% =============================================================================
% REVIEW_REPORT.tex --- Journal-reviewer report on MARST manuscript
% Compile: pdflatex REVIEW_REPORT.tex (twice for refs)
% =============================================================================

\documentclass[11pt]{article}

\usepackage[a4paper,left=1.1in,right=1.1in,top=1.0in,bottom=1.0in]{geometry}
\usepackage[utf8]{inputenc}
\usepackage[T1]{fontenc}
\usepackage{lmodern}
\usepackage{microtype}
\usepackage{textcomp}

\usepackage{booktabs}
\usepackage{array}
\usepackage{tabularx}
\usepackage{enumitem}
\setlist[itemize]{leftmargin=*,itemsep=2pt,topsep=2pt}
\setlist[enumerate]{leftmargin=*,itemsep=2pt,topsep=2pt}

\usepackage{titlesec}
\titleformat{\section}{\Large\bfseries\sffamily}{\thesection}{0.6em}{}
\titleformat{\subsection}{\large\bfseries\sffamily}{\thesubsection}{0.5em}{}

\usepackage{xcolor}
\definecolor{accent}{HTML}{1F4E79}
\definecolor{warn}{HTML}{B45F06}
\definecolor{ok}{HTML}{1B5E20}
\definecolor{boxbg}{HTML}{F5F7FA}
\definecolor{boxrule}{HTML}{B6C2CC}

\usepackage[most]{tcolorbox}
\tcbset{
  reviewerbox/.style={
    colback=boxbg, colframe=boxrule,
    boxrule=0.4pt, arc=2pt,
    left=8pt, right=8pt, top=4pt, bottom=4pt,
    fonttitle=\bfseries\sffamily\color{accent},
  }
}

\usepackage[colorlinks=true,
            linkcolor=accent, citecolor=accent, urlcolor=accent]{hyperref}

\usepackage{fancyhdr}
\pagestyle{fancy}
\fancyhf{}
\renewcommand{\headrulewidth}{0.4pt}
\fancyhead[L]{\small\sffamily MARST manuscript --- reviewer report}
\fancyhead[R]{\small\sffamily 2026-06-02}
\fancyfoot[C]{\thepage}

\title{\vspace{-1.2cm}\sffamily\bfseries
       MARST manuscript --- journal-reviewer report}
\author{}
\date{\vspace{-2.5em}\sffamily\small
      \textbf{Manuscript:} \texttt{ARTICLE.tex} (25 pages, 7{,}497 words,
      single-column NeurIPS/arXiv style)\\
      \textbf{Source notebook:} \texttt{marst-learnable-results.ipynb}\\
      \textbf{Date of review:} 2026-06-02\\
      \textbf{Reviewer role:} senior co-author preparing the paper for
      submission to \textbf{Transportation Research Part C} / \textbf{IEEE T-ITS}\\
      \textbf{Rubric:} novelty / technical correctness / experimental rigor /
      reproducibility / statistical validity / writing quality /
      publication readiness}

% =============================================================================
\begin{document}
\maketitle
\thispagestyle{fancy}

% =============================================================================
\section{Overall recommendation}

\noindent\textbf{Major revision needed before submission.}  The science
is sound, the leak audit and statistical machinery are unusually strong
for this venue, and the negative finding on per-sensor learnable decay
is genuinely creditable.  But three structural issues --- the
parameter-count argument, the absence of a load-balancing ablation, and
a sub-optimal KNN baseline --- will all be flagged by a competent
reviewer.  None requires re-running the main experiments; all can be
addressed in 1--2 days of focused work plus prose edits.

% =============================================================================
\section{Major strengths}

\begin{tabularx}{\linewidth}{@{}lX@{}}
\toprule
\textbf{\#} & \textbf{Strength} \\
\midrule
\textbf{S1} & \textbf{Strict streaming protocol with an empirical leak
              audit.}  The per-dataset held-out-truth corruption test
              plus future-perturbation test (\S3.5, \S5.8) is rare in
              this literature.  The 14-baseline
              \texttt{NO LEAK / CAUSAL} verdict is a hard claim that
              survives scrutiny and gives the paper a deployability
              story that BRITS-forward, iTransformer, etc.\ share but
              few discuss explicitly. \\[4pt]
\textbf{S2} & \textbf{Statistical rigor is appropriate for the venue.}
              3 training seeds $\times$ 5 eval-mask seeds = 15 paired
              evaluations, Wilcoxon signed-rank, Holm--Bonferroni across
              the four dataset comparisons, all four reach
              $p < 0.001$.  Table~3 cleanly separates the headline
              claim from noise. \\[4pt]
\textbf{S3} & \textbf{Honest negative finding.}  The per-sensor
              learnable decay parameter is reported as a near-null
              result (std $\approx 0.003$ around mean $\approx 0.95$)
              and the implication for the contribution claim is
              acknowledged in \S6.  Exactly the kind of self-critical
              discussion reviewers respect. \\[4pt]
\textbf{S4} & \textbf{JamMAE consistency check.}  The pipeline asserts
              that the MAE leader and JamMAE leader are the same model
              on every dataset.  Domain-appropriate operational measure
              rather than a synthetic add-on metric. \\[4pt]
\textbf{S5} & \textbf{Reproducibility is strong.}  Concrete
              hyperparameters, RNG discipline (per-thread NumPy and
              CUDA generators), 500-step train/eval gap, exact dataset
              slicing, runtime, and a named notebook --- Appendix~B is
              at the level a third party could reproduce the tables. \\
\bottomrule
\end{tabularx}

% =============================================================================
\section{Major weaknesses}

\subsection*{W1 --- No ablation on the load-balancing auxiliary loss}
The aux loss at $\lambda = 0.02$ is presented as a design move (\S3.4)
and as one of four contributions (\S1), but the manuscript shows no
$\lambda$-sweep, no $\lambda = 0$ comparison, and no quantitative
evidence that the load-balanced variant beats the unbalanced variant
on the headline metric.  The qualitative ``the gate collapses without
it, look at the ablation rows'' argument in \S5.7 is not a substitute
for a single comparison row $\lambda = 0$ vs $\lambda = 0.02$ on MAE.
\textbf{A reviewer will demand this.}
\textit{Effort to fix:} 4 hours on the same 2$\times$\,T4 setup.

\subsection*{W2 --- KNN baseline is sub-optimal --- and we have the data to prove it}
The Q12 KNN $k$-sweep in cell 15 (output:
\texttt{k=10 $\rightarrow$ MAE=19.60} on PEMS08, vs
\texttt{k=5 $\rightarrow$ MAE=20.18}) shows $k = 10$ is better on every
dataset, but the main results table reports $k = 5$.  A reviewer will
ask why we picked the worse setting.  Either re-run with $k = 10$ or
report the sweep explicitly and acknowledge the choice.  The current
paper makes neither move.
\textit{Effort:} 30 minutes.

\subsection*{W3 --- Parameter-count gap acknowledged but not addressed}
MARST is 17$\times$ iTransformer.  The Discussion promises a
compute-matched ablation as ``future work.''  A competent reviewer
(especially at T-ITS) will not accept this as sufficient --- they will
say ``run one 200\,K-parameter MARST seed and report the row.''
Without that row the win is vulnerable to the ``you bought it with
parameters'' rebuttal even if the architectural prior really is doing
work.
\textit{Effort:} 1 hour training + 15 min prose.

\subsection*{W4 --- ImputeFormer is cited but never compared, and the framing is now awkward}
After the recent edit, the paper says ``We do not run ImputeFormer as a
baseline \ldots it is not strictly causal.''  This is honest but
creates a problem: ImputeFormer is offline, so it is the strongest
published method on PEMS-BAY/METR-LA \emph{in the offline setting}.
Without an offline-vs-streaming comparison, a reviewer cannot tell
whether MARST is competitive with the offline state-of-the-art.  Two
clean options:
\begin{itemize}
\item add ImputeFormer's published numbers as a non-comparable
reference row in Table~2 with a footnote (``offline; not directly
comparable'');
\item remove the ImputeFormer mention entirely except as related
work.
\end{itemize}

\subsection*{W5 --- Missing-pattern table has only 3 of 4 datasets}
Table~7 reports PEMS-BAY, METR-LA, PEMS08; PEMS04 is absent.  Either
add it (the notebook prints it) or explain the omission.  Reviewers
spot this gap immediately.

\subsection*{W6 --- Per-sensor learnable decay is sold as a contribution and as a null result simultaneously}
Awkward positioning.  Recommend re-positioning: drop it from the
contributions list in \S1 (contribution \#2), describe it instead as a
\emph{design probe} in \S3.2 whose outcome --- a negative finding ---
is reported in \S5.5.  The Discussion already handles it well; the
Introduction should not promote it as a contribution.

\subsection*{W7 --- Only three sparsity levels (50, 80, 95\,\%)}
Standard in this benchmark is six (50, 60, 70, 80, 90, 95\,\%).
Defensible to keep three, but the omission of 70 and 90 should be
flagged in Limitations explicitly.  Currently it is mentioned in
passing.

% =============================================================================
\section{Required revisions (must fix before submission)}

\begin{tabularx}{\linewidth}{@{}lXl@{}}
\toprule
\textbf{\#} & \textbf{Action} & \textbf{Effort} \\
\midrule
\textbf{R1} & Add a single row to the results:
              \texttt{MARST ($\lambda = 0$)} vs
              \texttt{MARST ($\lambda = 0.02$)} MAE on all four
              datasets, same 3-seed protocol.  Without this row the
              load-balancing-loss contribution claim is unfounded.
              & 4 h \\[4pt]
\textbf{R2} & Either re-run the KNN imputer baseline with $k = 10$ (the
              Q12-sweep winner) or add a sentence in \S4 (``Baselines'')
              explaining the $k = 5$ choice and reporting the
              $k$-sweep result in a footnote.
              & 30 min \\[4pt]
\textbf{R3} & Add PEMS04 to the missing-pattern table (Table~7).  The
              data is in the notebook output.
              & 10 min \\[4pt]
\textbf{R4} & Move per-sensor learnable decay out of the contributions
              list in \S1.  Re-cast contribution \#2 as the
              load-balancing aux loss (more defensible); report
              per-sensor learnable decay as a design probe with a
              negative-finding outcome in \S3.2 / \S5.5.
              & 30 min \\[4pt]
\textbf{R5} & Add a row to Table~9 (parameter count) or a paragraph in
              \S6 with a compute-matched mini-MARST result (hidden =
              64, layers = 3, $\approx$\,200\,K parameters).  Even a
              single training seed is enough for the
              parameter-fairness rebuttal.
              & 1 h + 15 min \\
\bottomrule
\end{tabularx}

\medskip
\noindent\textbf{Total required-revision effort:} $\approx$\,1 day of
training compute on Kaggle 2$\times$\,T4 + 1.5 h of prose / table edits.

% =============================================================================
\section{Suggested revisions (would strengthen but not blocking)}

\begin{tabularx}{\linewidth}{@{}lX@{}}
\toprule
\textbf{\#} & \textbf{Action} \\
\midrule
\textbf{SR1} & Report JamMAE for every baseline in an extended-metrics
               table per dataset, not just LOCF.  The pipeline already
               computes it (cell 15 extended-metrics block). \\[3pt]
\textbf{SR2} & Add ImputeFormer's published numbers as a
               non-comparable reference row in Table~2 with a clear
               ``(offline)'' footnote so reviewers can see the offline
               state-of-the-art context.  Quoting published numbers is
               acceptable. \\[3pt]
\textbf{SR3} & Add one speed-dataset example to Figure~11 (currently
               only PEMS08 flow).  The qualitative panel is more
               informative with both regimes represented. \\[3pt]
\textbf{SR4} & Ablate the 2-hop neighbour-mean input feature
               ($n_{\text{2hop}}$).  Currently described as ``extra
               spatial context'' with no quantitative evidence.  One
               row is enough. \\[3pt]
\textbf{SR5} & Replace the present-tense
               \verb|\paragraph{Limitations.}| block with a numbered
               subsection (L1, L2, $\ldots$) so it is easier for
               reviewers to reference. \\[3pt]
\textbf{SR6} & Add 2025--2026 streaming-imputation references if any
               have appeared post-July 2025; current Related Work has
               BayOTIDE and STAMImputer but not much beyond. \\[3pt]
\textbf{SR7} & Reframe the ``soft-LOCF decay sweep'' in Table~6 as an
               \emph{initialisation sensitivity} test, not a
               hyperparameter sweep, to be precise about what is being
               measured. \\[3pt]
\textbf{SR8} & Section~6 has seven paragraphs; merge ``The headline win
               is real'' and ``Streaming deployment'' (both essentially
               recap \S5.1 framing), and split the over-long
               Limitations paragraph into its constituent points. \\
\bottomrule
\end{tabularx}

% =============================================================================
\section{Minor / cosmetic issues}

\begin{itemize}
\item \textbf{Figure 6} caption (\textsf{mae vs jammae}) shows PEMS08
only.  Consider a footnote that the same pattern holds on the other
three datasets.
\item \textbf{Table 1} (positioning matrix).  The ``Per-pos.\ mixture''
column has \textsf{yes (neural experts)} for STAMImputer; consider
adding \textsf{(classical anchors)} for MARST's row to make the
comparison crisp at a glance.
\item \textbf{\texttt{references.bib}} does not include Pi-Transformer;
required only if SR6 is adopted.
\item \textbf{Hyphenation.}  \textsf{JamMAE} is one word in the prose
but rendered as \textsf{Jam-MAE} or \textsf{JamMAE} inconsistently in
some captions.  Pick one.
\end{itemize}

% =============================================================================
\section{Simulated specialist reviews}

\begin{tcolorbox}[reviewerbox, title=Friendly reviewer (PhD supervisor tone)]
``Strong, deployable, honest paper.  The leak audit and JamMAE
consistency check are unusual and good.  Two things before submission:
(i)~the load-balancing ablation row, (ii)~the compute-matched
mini-MARST.  With those, ready to send.  Keep the negative finding on
per-sensor learnable decay --- that is a sign of a mature researcher,
not a weakness.''
\end{tcolorbox}

\begin{tcolorbox}[reviewerbox, title=Harsh reviewer (top conference tone)]
``The contribution surface is thin.  (1)~MARST is a mixture of three
textbook predictors with a transformer trunk and a load-balancing loss
adopted from Switch Transformer; nothing here is architecturally new.
(2)~The KNN baseline is run at $k = 5$ while the authors' own Q12
sweep shows $k = 10$ is better --- careless or selective.  (3)~The
parameter count is 17$\times$ the strongest baseline; the comparison
is not fair.  (4)~ImputeFormer is mentioned as state-of-the-art but
never compared.  (5)~The load-balancing claim is qualitative; show me
the $\lambda = 0$ row.  Reject without these fixes.  With them, weak
accept.''
\end{tcolorbox}

\begin{tcolorbox}[reviewerbox, title=TRC reviewer]
``The streaming constraint and leak audit address a real deployment
concern that the imputation community has under-discussed.  The
four-dataset benchmark and the statistical machinery are appropriate.
Three required revisions: (i)~PEMS04 in the missing-pattern table,
(ii)~the load-balancing aux-loss ablation, (iii)~clarification of the
KNN-$k$ choice.  The cross-domain validation gap is acceptable for
this venue; transportation-only is in scope.''
\end{tcolorbox}

\begin{tcolorbox}[reviewerbox, title=IEEE T-ITS reviewer]
``Methodology is sound, statistics are correct, reproducibility is
good.  The parameter-efficiency comparison is the weakest link --- I
want a compute-matched row.  The deployment claim (`strictly
streaming, leak-audited') is the differentiator and should be
foregrounded in the Introduction more prominently.  Also: report
JamMAE for every baseline in the extended-metrics table, not just
LOCF.''
\end{tcolorbox}

% =============================================================================
\section{Revision checklist (printable)}

\subsection*{Before submission --- required}
\begin{itemize}
\item[\fbox{\phantom{X}}] \textbf{R1.} Add $\lambda = 0$ vs
       $\lambda = 0.02$ ablation row to Results.
\item[\fbox{\phantom{X}}] \textbf{R2.} Switch KNN baseline to $k = 10$
       (or report sweep).
\item[\fbox{\phantom{X}}] \textbf{R3.} Add PEMS04 row to
       missing-pattern table.
\item[\fbox{\phantom{X}}] \textbf{R4.} Move per-sensor learnable decay
       out of Contributions list.
\item[\fbox{\phantom{X}}] \textbf{R5.} Add compute-matched mini-MARST
       row (hidden 64, 3 layers).
\end{itemize}

\subsection*{Suggested --- would strengthen}
\begin{itemize}
\item[\fbox{\phantom{X}}] \textbf{SR1.} Report JamMAE for every baseline.
\item[\fbox{\phantom{X}}] \textbf{SR2.} ImputeFormer reference row in Table~2.
\item[\fbox{\phantom{X}}] \textbf{SR3.} Speed-dataset example in qualitative figure.
\item[\fbox{\phantom{X}}] \textbf{SR4.} 2-hop neighbour-mean feature ablation.
\item[\fbox{\phantom{X}}] \textbf{SR5.} Numbered limitations subsection.
\item[\fbox{\phantom{X}}] \textbf{SR6.} Newer streaming-imputation citations.
\item[\fbox{\phantom{X}}] \textbf{SR7.} Reframe decay-sweep table as initialisation sensitivity.
\item[\fbox{\phantom{X}}] \textbf{SR8.} Tighten Discussion section structure.
\end{itemize}

\subsection*{Cosmetic}
\begin{itemize}
\item[\fbox{\phantom{X}}] Figure 6 caption --- note cross-dataset generalisation.
\item[\fbox{\phantom{X}}] Table 1 --- add ``(classical anchors)'' to MARST's row.
\item[\fbox{\phantom{X}}] Hyphenation pass on \texttt{JamMAE}.
\end{itemize}

% =============================================================================
\section{Bottom line}

This is a \textbf{publishable paper after R1--R5 are addressed.}  None
of the required revisions invalidate the existing results; they fill
gaps that any reviewer will spot.  The science as it stands is honest,
the negative finding on per-sensor learnable decay is genuinely
creditable, and the streaming + leak-audit story is the strongest
single asset.

R3 and R4 can be executed in the current session without new
experiments.  R1, R2, R5 require notebook re-runs on Kaggle.

\bigskip\noindent
\textit{Generated 2026-06-02 from \texttt{ARTICLE.tex} and
\texttt{marst-learnable-results.ipynb}.}

\end{document}
